%-------------------------------------------------------------------------------
% SECTION TITLE
%-------------------------------------------------------------------------------
\cvsection{Expérience}

%-------------------------------------------------------------------------------
% CONTENT
%-------------------------------------------------------------------------------
\begin{cventries}

\cventry
{Assistant de recherche (Temps partiel permanent)} % Job title
{IMC2 -- Institut multidisciplinaire en cybersécurité et cyberrésilience} % Organisation
{Distance (Montréal, QC)} % Location
{Mar 2025 -- Fév 2026} % Date(s)
{
\begin{cvitems}
\item {Piloté GitLab CI/CD et runners pour services de recherche sur GCP ; charges Docker, exploitation du réseau GKE et livraisons d'infrastructure avec Terraform, \uline{renforçant des mises en production cloud reproductibles}}
\item {Renforcé microservices Java avec JUnit et Playwright ; promotion des builds Android via GitLab vers le GKE géré par Terraform, \uline{haussant la couverture avant release et alignant mobile sur le backend}}
\item {Livré une analyse automatisée de risques cybersécurité sur \uline{1000+ risques} et \uline{181+ catégories de menaces}, avec score pondéré, priorisation par étoiles et interface Angular de triage pour \uline{un triage analyste plus rapide et homogène}}
\item {\textbf{Compétences techniques :} Java, Kotlin, Android, Docker, Kubernetes, Terraform, GitLab CI/CD, GCP, JUnit, Playwright, Angular}
\end{cvitems}
}

\cventry
{Chargé de travaux pratiques, LOG8100 DevSecOps (Contrat temps partiel)} % Job title
{Polytechnique Montréal} % Organisation
{Montréal, QC (Sur place)} % Location
{Sep 2024 -- Jan 2026} % Date(s)
{
\begin{cvitems}
\item {Structuré et maintenu l'espace GitLab (dépôts, modèles CI/CD, runners, workflows) pour \uline{\textasciitilde20 équipes en parallèle}, \uline{réduisant la friction administrative}}
\item {Coaching Docker et Kubernetes (construction, déploiement, débogage) selon LOG8100 ; exemples sécurité inspirés d'OWASP, \uline{reliant théorie à des scénarios concrets}}
\item {Actualisé travaux et supports DevSecOps ; correction de TP et coaching sur remises, \uline{améliorant qualité et délais de correction}}
\item {\textbf{Compétences techniques :} GitLab, CI/CD, DevSecOps, Docker, Kubernetes, OWASP, codage sécurisé, Linux}
\end{cvitems}
}

\cventry
{Chargé de travaux pratiques (Contrat temps partiel)} % Job title
{École de technologie supérieure (ÉTS)} % Organisation
{Montréal, QC (Hybride)} % Location
{Fév 2024 -- Avr 2026} % Date(s)
{
\begin{cvitems}
\item {\uline{Environ 150 heures de contact} en laboratoires, soutien intégrateur et permanences sur parcours mobile/UX, intégrateur et données distribuées (p.\,ex. TCH057, IND500, TCH099, GTI660, GTI320)}
\item {Conçu et mis à jour des exercices SQL, MongoDB, Kafka et PostgreSQL avec grilles pour IND500, GTI660 et cohortes associées, \uline{alignant l'évaluation sur les résultats d'apprentissage}}
\item {Débloqué équipes intégrateur sur architecture et pipelines backend/données ; maintenu l'infrastructure de cours sur Azure et le coaching SCRUM ; animé mobile/UX et GTI320, \uline{pour que les étudiants livrent des parcours bout-en-bout}}
\item {\textbf{Compétences techniques :} SQL, MongoDB, Kafka, PostgreSQL, bases de données distribuées, Microsoft Azure, Android Studio, UI/UX, SCRUM}
\end{cvitems}
}

\cventry
{Développeur cloud (Temps plein permanent)} % Job title
{IONODES} % Organisation
{Laval, QC (Hybride)} % Location
{Mai 2023 -- Août 2023} % Date(s)
{
\begin{cvitems}
\item {Conçu un \uline{modèle d'abonnement Auth0 à trois niveaux} pour l'accès IoT multi-organisation, segmentant les capacités par locataire}
\item {Intégré Sentry en C\# avec Serilog pour erreurs structurées en production, \uline{accélérant le diagnostic pour support et développement}}
\item {Renforcé la préparation streaming ONVIF et WebRTC à l'échelle IoT ; orchestré Azure DevOps avec feature flags et promotions reproductibles, \uline{réduisant le risque de mise en service}}
\item {\textbf{Compétences techniques :} C\#, .NET, Azure DevOps, Auth0, Sentry, Serilog, UML, microservices, ONVIF, WebRTC}
\end{cvitems}
}

\cventry
{Développeur Back End (Coopérative)} % Job title
{Intact} % Organisation
{Montréal, QC (Hybride)} % Location
{Jan 2022 -- Mai 2022} % Date(s)
{
\begin{cvitems}
\item {Conçu des campagnes d'intégration sensibles aux dates d'inscription et des flux API d'enrôlement sur systèmes hérités, \uline{accélérant l'activation des couvertures}}
\item {Livré des API Kotlin de connexion et MFA pour applis mobiles client tout en respectant les politiques}
\item {Enrichi la pile ELK avec tableaux Kibana et codes d'erreur harmonisés entre microservices, \uline{raccourcissant le triage incident pour les opérations}}
\item {\textbf{Compétences techniques :} Kotlin, Spring Boot, GraphQL, MongoDB, microservices, Docker Swarm, Jira, Kibana, SAFe}
\end{cvitems}
}

\cventry
{Spécialiste Web (Contrat temps partiel)} % Job title
{Divine Essence (Union-Nature)} % Organisation
{Dorval, QC (Sur place)} % Location
{Juin 2020 -- Déc 2020} % Date(s)
{
\begin{cvitems}
\item {Modernisé le site commerce bilingue (navigation blog FR/EN, accessibilité) et étiquettes de vente CommerceBuild orientées SEO, \uline{améliorant l'UX et le trafic qualifié vers les fiches produit}}
\item {Implémenté un flux API «\,Ajouter au panier\,» sur CommerceBuild ; corrigé régressions pop-up, bannière et grille ; intégré un chatbot 3CX, \uline{fluidifiant le checkout et le support en libre-service}}
\item {\textbf{Compétences techniques :} HTML, CSS, JavaScript, CommerceBuild, 3CX, SEO}
\end{cvitems}
}

\cventry
{Développeur Web Full Stack (Stages)} % Job title
{Power Go} % Organisation
{Montréal, QC} % Location
{Sep 2020 -- Avr 2021} % Date(s)
{
\begin{cvitems}
\item {Livré un analyseur Python/Django/Vue.js/PostgreSQL unifiant les flux fournisseurs XML, JSON, CSV et Excel pour l'inventaire, \uline{remplaçant des scripts ponctuels}}
\item {Intégré les API Facebook Shop et Marketplace pour la synchro catalogue ; service de taux CAD/USD lié à l'inventaire et aux dates d'effet, \uline{évitant la double saisie et alignant les catalogues transfrontaliers}}
\item {\textbf{Compétences techniques :} Python, Django, Vue.js, PostgreSQL, MongoDB, JavaScript, Git, Jira, développement full-stack}
\end{cvitems}
}

\cventry
{Développeur Full Stack (Stages $\rightarrow$ Temps partiel)} % Job title
{Wandrian} % Organisation
{Montréal, QC} % Location
{Mai 2019 -- Janvier 2020} % Date(s)
{
\begin{cvitems}
\item {Sur la billetterie Italiarail avec Railkey Tech : performance et UX, Braintree, UI alignée à la marque et APIs d'achat, \uline{stabilisant les réservations bout-en-bout}}
\item {Réécrit et étendu la suite Selenium E2E ; outils internes Django, Pyramid et jQuery et tableau d'usage des devises, \uline{interceptant des régressions UI tôt et remplaçant des feuilles ad hoc pour CS/finance}}
\item {\textbf{Compétences techniques :} Python, Django, Pyramid, JavaScript, jQuery, Selenium, MongoDB, Git, Jira, automatisation de tests}
\end{cvitems}
}

\cventry
{Stagiaire en ingénierie} % Job title
{Algolux} % Organisation
{Montréal, QC (Sur place)} % Location
{Mai 2017 -- Août 2017; Mai 2018 -- Août 2018} % Date(s)
{
\begin{cvitems}
\item {Automatisé la validation caméra/ISP avec Ansible, VM Linux/QEMU UEFI, imagerie HDR Nexus~6 et services REST Flask pour capture et affichage à distance, \uline{réduisant les boucles manuelles}}
\item {Produit des jeux de données vision extérieure via campagne capture/annotation \uline{\textasciitilde12k images}, outillage Python couleur RAW et commande de rigs ; soutien postes QNAP, réseau et serveurs Linux/Windows pour les expériences}
\item {\textbf{Compétences techniques :} Python, C++, CUDA, OpenCV, Flask, Ansible, Linux, Windows}
\end{cvitems}
}

\end{cventries}
